\chapter{Use Case: EPOS API}
\label{chap:caso-epos}

This chapter applies the framework of Chapter~\ref{chap:metodo-arq} to the
production \ac{API}, \texttt{api\_v4.0}. The \ac{EPOS} itself~\cite{epos} is
a pan-European research infrastructure in solid Earth science, organised
into Thematic Core Services for seismology, geomagnetic observations, \ac{GNSS}
data and products~\cite{eposGnss}, and so on; the \ac{SEGAL} laboratory of
\ac{UBI} contributes to the \ac{GNSS} service by operating a node that
hosts station metadata and \ac{RINEX}~\cite{rinex} observation files, and
\texttt{api\_v4.0} is the public interface of that node. Three properties
of the domain shape the test framework: the entities are relatively static (stable
baselines are easy to obtain), the typical client request returns a long
list of records (the comparison algorithm must scale to long diffs), and the
same information is exposed in several serialization formats (the surface
area to test multiplies but the underlying logic does not).

\section{Architecture of \texttt{api\_v4.0}}
\label{sec:apiv4-arch}

The application is a Maven \ac{WAR} project targeting Java SE 21 and Jakarta
EE 11, deployed onto Payara Server 6 with \texttt{finalName=ROOT} so that it
serves the root context. Its source tree is organised into four top-level
packages that form a classical layered architecture (Figure~\ref{fig:apiv4-layers}):
\texttt{Glass} provides the \ac{REST} layer (\ac{JAX-RS} resources \texttt{StationEndpoints},
\texttt{FilesEndpoints}, \texttt{ListEndpoints} and \texttt{Start}, plus
parsers, format writers and the validation/shape helpers); \texttt{Geometry}
and the helper classes \texttt{Validate.java} and \texttt{Shape.java} form
a thin business layer; \texttt{DB\_Tables} holds the forty \ac{JPA} entities
that model the relational schema (stations, networks, agencies, file types,
observation summaries, embargoes); and \texttt{Config} (\texttt{Loader.java})
is the cross-cutting concern that handles bootstrap.

Solid arrows denote ordinary call/data flow; dashed arrows denote
configuration and infrastructure binding established at deployment time.

\begin{figure}[h]
	\centering
	\resizebox{\linewidth}{!}{%
	\begin{tikzpicture}[
		every node/.style={font=\small},
		node distance=4mm,
		layer/.style={draw, rounded corners=2pt, align=center, thick, minimum width=92mm, minimum height=12mm},
		cc/.style={draw, dashed, rounded corners=2pt, align=center, minimum width=28mm, minimum height=10mm, font=\footnotesize},
		arrow/.style={-{Stealth[length=2mm]}, thick}
	]
		\node[layer, fill=dkblue!10]
			(rest)
			{\textbf{REST layer}, package \texttt{Glass} \\ \footnotesize Start, StationEndpoints, FilesEndpoints, ListEndpoints, Format writers};
		\node[layer, fill=mauve!10, below=of rest]
			(biz)
			{\textbf{Business / validation layer} \\ \footnotesize \texttt{Validate.java}, \texttt{Shape.java}, package \texttt{Geometry} (\texttt{NativePolygonCheck})};
		\node[layer, fill=dkgreen!10, below=of biz]
			(jpa)
			{\textbf{Persistence layer}, package \texttt{DB\_Tables} \\ \footnotesize 40 JPA entities (Station, RinexFile, Network, Agency, \ldots)};
		\node[layer, fill=dark!10, below=of jpa]
			(db)
			{\textbf{Database}, PostgreSQL \texttt{192.168.168.30:5432/DB\_3.0\_J} \\ \footnotesize over the SEGAL internal network / \ac{VPN}};

		\node[cc, right=6mm of biz, yshift=8mm]
			(cfg)
			{\textbf{Config} \\\texttt{Loader.java} \\ env / file / default};
		\node[cc, right=6mm of biz, yshift=-8mm]
			(pu)
			{Persistence unit \\\texttt{MeuPU} \\ HikariCP pool};

		\draw[arrow] (rest) -- (biz);
		\draw[arrow] (biz) -- (jpa);
		\draw[arrow] (jpa) -- (db);
		\draw[arrow, dashed] (cfg) -- (db);
		\draw[arrow, dashed] (pu) -- (jpa);
	\end{tikzpicture}%
	}
	\caption{Layered architecture of \texttt{EPOS V4}.}
	\label{fig:apiv4-layers}
\end{figure}

\subsection{Endpoints and Output Formats}

The \ac{REST} surface is declared in four classes inside the \texttt{Glass}
package. \texttt{Start.java} carries the \texttt{@ApplicationPath("/api")} annotation
and registers a \texttt{PingResource} for the trivial \texttt{/api/ping} liveness
check. The three substantive resources are \texttt{StationEndpoints} (\texttt{/api/stations},
seven endpoints for station metadata indexed by identifier, location,
agency, network, equipment or coordinates), \texttt{FilesEndpoints} (\texttt{/api/files},
four endpoints for \ac{RINEX} file metadata filtered by agency/antenna/bounding
box/date, by marker, the high-rate variant, and an aggregated combination),
and \texttt{ListEndpoints} (\texttt{/api/list}, a parametric endpoint
returning the valid values of a given filter). Every endpoint accepts a
\texttt{\{format\}} path parameter and emits \ac{JSON}, \ac{XML}, \ac{CSV}
or (for the file endpoints) a downloadable shell script; the same \ac{OpenAPI}
3.1.0 document is served at \texttt{/swagger.html} through a Swagger \ac{UI}
hosted inside the \ac{WAR}.

\subsection{Persistence}

The persistence unit, declared in \texttt{src/\allowbreak main/\allowbreak
resources/\allowbreak META-INF/\allowbreak persistence.xml}, is named
\texttt{MeuPU} (inherited from the legacy code). It is configured as
\texttt{RESOURCE\_LOCAL}, with Hibernate as the \ac{JPA} provider and a connection
pool managed by HikariCP\@. Database connection parameters, host, port,
database name, user, password, are resolved at startup by \texttt{Config/Loader.java},
with the following precedence: environment variables, then a \texttt{Config.config}
file at the root of the deployment, then hard-coded defaults pointing at the
laboratory's database (\texttt{192.168.168.30:5432}).

The laboratory maintains two database fixtures with the same schema but
different volumes of data:

\begin{itemize}
	\item \texttt{DB\_3.0\_J} hosts a small but representative production
		subset, on the order of ten stations, eighty-six networks, two hundred and
		thirty agencies and roughly one hundred and forty-six thousand \texttt{rinex\_file}
		rows at the time of writing. This is the fixture used by the correctness
		diff runs, because it is small enough to allow per-case comparison of
		full response bodies without overwhelming the report.

	\item \texttt{DB\_3.0\_F} hosts the full dataset and is used by the
		volume-oriented diff runs, where the goal is not byte-for-byte equality but
		the detection of changes that would only be apparent at scale, for
		example, an inadvertent \texttt{N+1} query introduced by a refactor, or a
		query plan regression visible only on a large table.
\end{itemize}

The precedence rules of \texttt{Loader.java} are what allow the \texttt{http-diff-job}
to switch the running \acp{WAR} between the two fixtures without rebuilding
the artifact: the job sets the appropriate environment variables in its \texttt{variables:}
block and the application picks them up at boot.

The shape of the schema is discussed in Section~\ref{sec:epos-db}; a simplified
view of the central entities is given in Figure~\ref{fig:datamodel}.

Representative columns as they appear in the live PostgreSQL schema. Green
entities carry domain data; purple entities are controlled-vocabulary lookups
(note that \texttt{country.iso\_code} is the natural primary key, not a
surrogate \texttt{id}); grey entities are junction tables that materialise the
many-to-many relationships of the model. Arrows point from the ``one'' side
to the ``many'' side of the relationship.

\begin{figure}[h]
	\centering
	\resizebox{\linewidth}{!}{%
	\begin{tikzpicture}[
		every node/.style={font=\footnotesize},
		entity/.style={draw, rounded corners=2pt, align=center, thick, minimum width=24mm, minimum height=11mm, fill=dkgreen!10},
		lookup/.style={entity, fill=mauve!10},
		junc/.style={entity, fill=lgray, minimum width=22mm},
		arrow/.style={-{Stealth[length=2mm]}, thick},
		rel/.style={font=\scriptsize, text=dark}
	]
		\node[entity]
			(station)
			{\textbf{station} \\\scriptsize id, marker[9], lat, lon, x/y/z,\\\scriptsize country\_code, station\_type};
		\node[junc, left=10mm of station]
			(sn)
			{\textbf{station\_network} \\\scriptsize include\_date};
		\node[entity, left=10mm of sn]
			(network)
			{\textbf{network} \\\scriptsize id, name, network\_type, doi};
		\node[junc, left=10mm of network]
			(an)
			{\textbf{agency\_network} \\\scriptsize role};
		\node[lookup, left=10mm of an]
			(agency)
			{\textbf{agency} \\\scriptsize id, name, abbreviation,\\\scriptsize ror};
		\node[lookup, above=8mm of station]
			(country)
			{\textbf{country} \\\scriptsize iso\_code (PK), name};
		\node[entity, right=12mm of station]
			(rinex)
			{\textbf{rinex\_file} \\\scriptsize id, name, file\_size,\\\scriptsize reference\_date, md5checksum};
		\node[entity, below=8mm of rinex]
			(highrate)
			{\textbf{highrate\_rinex\_file}};
		\node[lookup, right=12mm of rinex]
			(ftype)
			{\textbf{file\_type} \\\scriptsize format, sampling\_window};
		\node[junc, below=8mm of station] (nodest) {\textbf{node\_station}};
		\node[entity, below=8mm of nodest]
			(node)
			{\textbf{node} / \textbf{data\_center}};

		\draw[arrow] (agency) -- node[rel, above] {1..*} (an);
		\draw[arrow] (network) -- node[rel, above] {1..*} (an);
		\draw[arrow] (station) -- node[rel, above] {1..*} (sn);
		\draw[arrow] (network) -- node[rel, above] {1..*} (sn);
		\draw[arrow] (country) -- node[rel, right] {1..*} (station);
		\draw[arrow] (station) -- node[rel, above] {1..*} (rinex);
		\draw[arrow] (station) -- node[rel, above, sloped] {1..*} (highrate);
		\draw[arrow] (ftype) -- node[rel, above] {1..*} (rinex);
		\draw[arrow] (ftype) -- node[rel, below, sloped] {1..*} (highrate);
		\draw[arrow] (station) -- node[rel, right] {1..*} (nodest);
		\draw[arrow] (nodest) -- node[rel, right] {*..1} (node);
	\end{tikzpicture}%
	}
	\caption{Simplified view of the core \ac{EPOS} data model}
	\label{fig:datamodel}
\end{figure}

\section{The CI/CD Pipeline}
\label{sec:pipeline}

The pipeline of \texttt{api\_v4.0} extends the patterns validated on the first
prototype with an additional job in a new \texttt{diff} stage. The diff job compares
the current build against \emph{two} baselines, in what amounts to a three-way
comparison: the immediate parent commit \texttt{HEAD\textasciitilde 1} (which
catches regressions introduced by the commit currently being pushed) and a \emph{fixed}
reference commit identified by a SHA stored in the \texttt{DIFF\_FIXED\_BASELINE}
GitLab \ac{CI}/\ac{CD} variable (which catches the slow drift that
accumulates over many commits along a long-running refactor branch). The
key parts of the \texttt{.gitlab-ci.yml} are shown, condensed, in Listing~\ref{app:glass-ci-cd}.

Three details of the listing are worth highlighting. The \emph{recent} baseline
is \texttt{HEAD\textasciitilde 1}, which keeps the per-commit comparison focused
on the change just introduced; the \emph{fixed} baseline is a commit \ac{ID}
stored in \texttt{DIFF\_FIXED\_BASELINE} (typically pointing at the last
commit known to have shipped to the \ac{EPOS} portal) and acts as a long-term
reference against which the cumulative drift of the branch is measured. Each
baseline build is materialised through \texttt{git worktree} rather than
\texttt{git checkout}, so the working copy is never disturbed and the
harness can freely compare files from both checkouts in parallel. Finally,
each baseline is optional and treated independently: if \texttt{HEAD\textasciitilde
1} does not exist (first commit of the repository) or \texttt{DIFF\_FIXED\_BASELINE}
is unset (start of a refactor branch), the corresponding sub-run is skipped.
A pipeline can therefore have zero, one or two diff sub-runs depending on
the situation, and never fails because a baseline is simply unavailable.

\section{The HTTP-Level Diff: \texttt{HttpDiff.java}}
\label{sec:httpdiff}

\texttt{HttpDiff.java} is the harness that exercises the \ac{API} end-to-end.
Its execution proceeds in four phases, summarised graphically in Figure~\ref{fig:httpdiff-phases}.
The \ac{HTTP} harness runs once and internally boots \emph{all} the available
\acp{WAR} in parallel (the \texttt{HEAD} \ac{WAR}, the recent baseline if
any, and the fixed baseline if any). It then issues every test case against every
running instance and emits one diff report per baseline.

The harness boots one Payara Micro per available \ac{WAR}, the \texttt{HEAD}
build and whichever of the recent/fixed baselines were provided, waits for all
of them to be ready, fires every test case against every instance in
parallel, then performs a three-way comparison and emits one diff report per
baseline.

\begin{figure}[h]
	\centering
	\resizebox{\linewidth}{!}{%
	\begin{tikzpicture}[
		every node/.style={font=\footnotesize},
		proc/.style={draw, rounded corners=2pt, align=center, thick, minimum width=26mm, minimum height=10mm, fill=dkblue!10},
		micro/.style={proc, fill=mauve!10, minimum width=30mm},
		head/.style={proc, fill=dkgreen!10, minimum width=30mm},
		stage/.style={font=\small\bfseries, text=dark},
		barrow/.style={-{Stealth[length=2mm]}, thick},
		darrow/.style={-{Stealth[length=2mm]}, thick, dashed}
	]
		\node[proc]
			(harness)
			{\textbf{HttpDiff.java} \\ \footnotesize (JUnit 5)};
		\node[head, right=18mm of harness]
			(p80)
			{\textbf{payara-micro :8080} \\ \footnotesize HEAD WAR};
		\node[micro, right=8mm of p80]
			(p81)
			{\textbf{payara-micro :8081} \\ \footnotesize recent baseline};
		\node[micro, right=8mm of p81]
			(p82)
			{\textbf{payara-micro :8082} \\ \footnotesize fixed baseline};
		\node[proc, fill=dark!10, right=12mm of p82]
			(pg)
			{\textbf{PostgreSQL} \\ \footnotesize 192.168.168.30};

		\node[stage, below=4mm of harness] (s1) {Phase 1, Boot};
		\draw[barrow] ($(harness.south)+(0,-12mm)$) -| (p80.south);
		\draw[barrow] ($(harness.south)+(0,-12mm)$) -| (p81.south);
		\draw[barrow] ($(harness.south)+(0,-12mm)$) -| (p82.south);

		\node[stage, below=20mm of s1] (s2) {Phase 2, Verify};
		\draw[barrow]
			($(harness.south)+(0,-30mm)$) -- ++(1.8,0)
			node[midway, above, font=\scriptsize] {GET /api/ping}
			-|
			(p80);
		\draw[barrow]
			($(harness.south)+(0,-34mm)$) -- ++(4.6,0)
			node[midway, above, font=\scriptsize] {GET /api/ping}
			-|
			(p81);
		\draw[barrow]
			($(harness.south)+(0,-38mm)$) -- ++(7.4,0)
			node[midway, above, font=\scriptsize] {GET /api/ping}
			-|
			(p82);

		\node[stage, below=24mm of s2] (s3) {Phase 3, Replay \& 3-way diff};
		\draw[barrow]
			($(harness.south)+(0,-56mm)$) -- ++(1.8,0)
			node[midway, above, font=\scriptsize] {case[i]}
			-|
			(p80);
		\draw[barrow]
			($(harness.south)+(0,-60mm)$) -- ++(4.6,0)
			node[midway, above, font=\scriptsize] {case[i]}
			-|
			(p81);
		\draw[barrow]
			($(harness.south)+(0,-64mm)$) -- ++(7.4,0)
			node[midway, above, font=\scriptsize] {case[i]}
			-|
			(p82);
		\draw[darrow]
			($(p80.south)+(0,-12mm)$) -- ($(pg.south)+(0,-12mm)$)
			node[midway, below, font=\scriptsize] {all three query the same DB};

		\node[
			proc,
			fill=dkgreen!10,
			below=70mm of harness,
			minimum width=132mm
		]
			(diff)
			{\textbf{Diff} HEAD vs recent $\rightarrow$ \texttt{http-diff-report/recent/} \quad \textbf{Diff} HEAD vs fixed $\rightarrow$ \texttt{http-diff-report/fixed/}};

		\node[stage, below=4mm of diff] (s4) {Phase 4, Teardown + Report};
		\draw[barrow]
			($(harness.south)+(0,-92mm)$) -- ++(10.0,0)
			node[midway, above, font=\scriptsize] {SIGTERM all payara-micro};
	\end{tikzpicture}%
	}
	\caption{Phases of \texttt{HttpDiff.java}.}
	\label{fig:httpdiff-phases}
\end{figure}

The figure already covers the gist; a few implementation details are worth recording.
The test cases live in \texttt{src/\allowbreak test/\allowbreak resources/\allowbreak
http-diff-cases.json} as a \ac{JSON} array of request descriptors (method,
path, headers, optional body); This guards against spurious failures from
requests landing on still-deploying containers. Bodies are compared as normalised
strings (sorted \ac{JSON} keys, trimmed whitespace) and the headers that vary
by definition (\texttt{Date}, \texttt{Server}, \texttt{Content-Length}) are
ignored. The two comparisons (HEAD vs recent, HEAD vs fixed) are kept in
separate report directories so that a per-commit regression and an
accumulated drift remain distinguishable in the artifact tree.

\section{Reports and Findings}
\label{sec:reports}

Each pipeline run publishes the integration-test reports, the diff reports
(\texttt{recent/} and \texttt{fixed/} subdirectories of \texttt{target/http-diff-report/})
and the JUnit XML files consumed by GitLab. Integration test results expire
after one week; the diff reports are kept for four because they are the historical
record of behavioural change across the refactor. During development the
harness surfaced real regressions hidden inside seemingly cosmetic
refactors that no line-by-line review of the diff would plausibly have
caught by eye. The next chapter draws the main conclusions of the project and
outlines the items that, by choice or by constraint, were left out of its scope.